% 08_discussion.tex — Discussion. Drafted by item P2-E8-9 (P2 Sprint 7).
% Voice: ../STYLE.md (Andreas Maier) — three-beat (claim -> Yet, -> Hence,),
% short paratactic declaratives, anaphoric recap to open, honest limitations,
% cross-disciplinary reach, close on a forward-looking programme (not a result).
% Numbers: tools/xai_study/compare/out/leaderboard.md + faithful_demo.md (committed).
% Honesty contract (../plan.md §"What we claim about semantics"): we DEFINE and
% MEASURE the semantic gap; we do NOT claim any method recovered meaning, algorithm,
% variables, or design from scratch. All T3 semantic labels are external
% (AtariARI/OCAtari) and only causally VERIFIED, not discovered. The plausibility
% axis is a documented method-tradition PROXY, not a human-subjects measurement.

\section{Discussion}\label{sec:discussion}

In this Article, we have turned the interpretability toolkit onto a machine whose
ground truth is computable, and we could measure two things at once. We have shown,
first, that on a system where the true cause of any output is recoverable by exact
intervention, causal and intervention methods are faithful while the field's most
familiar gradient and correlational tools are not---the leaderboard separates them
along the only axis that admits an exact answer (Fig.~\ref{fig:faithfulness}). We
have shown, second, that this is the easier half of the problem. Even a method that
scores at the oracle ceiling returns the causal \emph{wiring}---an effect map, a
data-flow graph---and not the \emph{semantics}, the variable's meaning or the
routine's algorithm. The residual between the two is not a matter of taste. On this
machine it is a number, and we have named it the \emph{semantic gap}. We discuss
here what that gap implies for interpretability practice, why a 1977 chip is a fair
place to measure it, where our claim stops, and what would be needed to close it.

\subsection*{Many traditions, one ground truth --- and the boundary it cannot
cross}

The first lesson is a confirmation, made quantitative. Three research traditions
that rarely share a page---a neuroscience battery, the attribution toolkit, and
mechanistic interpretability---turn out to occupy a single plane once a common
oracle scores them (Fig.~\ref{fig:faithfulness}). Jonas and Kording warned that an
analyst armed with the methods of systems neuroscience could fail to understand even
a microprocessor, because finding structure is not the same as explaining
it~\cite{jonas2017could}. We have made their lesson a measurement: the classical
battery recovers rich structure yet scores low against the known mechanism, the
gradient family collapses to the floor on the discrete position regime, and only the
causal operators reach the ceiling. Mechanistic interpretability, in particular,
has until now been calibrated against \emph{synthetic} or \emph{compiled} ground
truth---programs compiled into transformers by construction, or semi-synthetic
networks with a planted circuit~\cite{lindner2023tracr,gupta2024interpbench}; ground-truth
\emph{attribution} benchmarks likewise plant the important features~\cite{yang2019bim}.
We add the missing case: a real, hand-authored, deployed artifact, never designed to
be interpretable, whose ground truth is \emph{independent} of the artifact rather
than built into it. \emph{Hence} the leaderboard is not only a ranking; it is the
first ground-truth calibration of these methods on a system the field did not get to
plant.

Yet the same oracle marks a boundary the leaderboard cannot cross, and that boundary
is the real subject of this paper. A perfect faithfulness score certifies that a
method named the true causes. It does not certify that the method recovered what the
program \emph{means}. We saw the boundary three times, from three sides
(Fig.~\ref{fig:taxonomy}): a circuit that is graph-correct and minimal
($F=M=1.0$) yet reproduces behavior only $44\%$ of the time ($S=0.44$); a feature
dictionary that matches every named variable (matched fraction $1.0$) while being
nearly causally inert ($F=0.041$, $S=-0.336$); and a concept that is linearly
decodable yet, by the oracle, never used in the computation. The gap runs in both
directions on the same axis at once---correct-but-incomplete and named-but-wrong---
which is precisely what a ``methods are not good enough yet'' story cannot produce.

\subsection*{What \emph{understanding} would require --- the semantic level}

Why does faithfulness stop short? The cleanest answer comes from Marr's levels of
analysis~\cite{marr1982vision}. A faithful attribution lives at the implementational
level: it tells us which cells and registers carry the signal and how a perturbation
propagates. Understanding lives one level up, at the algorithmic level---what
function the program computes and how. A correct data-flow graph is an
implementational fact; the bounce algorithm it implements is an algorithmic one, and
the two are objects of a different type. This is also Lazebnik's lament, that one can
characterize a system in exhaustive mechanistic detail and still not understand
it~\cite{lazebnik2002can}. Our contribution is to turn that lament into a measured
quantity: the distance from a faithful effect table to the algorithm is the
sufficiency gap $1-S$, and on the exact effect table that gap cannot be closed by a
better hyperparameter, because the table is \emph{already} exact.

The sharpest way to state the bar is to borrow it from software engineering.
IEEE~610.12 defines software as programs \emph{plus} their documentation---the
specification, the design, the intent~\cite{ieee610}. A faithful circuit recovers
part of the program. Understanding the system means recovering its documentation:
what each variable is \emph{for}, what each routine is \emph{meant} to do, the
design the engineer held in mind. Our causal oracle, by construction, cannot reach
that object, because the meaning of a variable is not a causal fact about the
variable---it is a fact about the design that the variable serves. \emph{Hence} the
gap we measure is not an artifact of weak methods but a property of the question:
faithfulness is necessary for understanding, and it is not sufficient. We are
deliberate that this paper \emph{defines and measures} the gap. \emph{Closing} it
---recovering the program's documentation and design from the bare machine---is the
reverse-engineering bar, and it is the companion design-recovery paper (Paper~4),
which frames itself as measuring the distance to that bar, not clearing it. The
behavioral-grounding bridge---asking whether a recovered concept matches how the
system is \emph{used}---is the companion behavior paper (Paper~3). The learned agent,
where there is no disassembly to check against, is Paper~5.

A second strand of the literature reads our result the other way around. Probing
shows that a concept is \emph{present} in the state; it does not show that the
concept is \emph{used}, the control-task warning of Hewitt and
Liang~\cite{hewitt2019designing}. Our intervention oracle is strictly stronger: it
asks whether the byte \emph{causes} the output, and it found a decodable score cell
that the program never reads. \emph{Note that} this is the same boundary seen from
the data side. Decodability is an implementational presence; use is an algorithmic
fact; understanding needs the second, and only the oracle supplies it.

\subsection*{The shared toolkit, read across disciplines}

It is worth saying plainly that the three traditions are, at bottom, one toolkit
under different names. Tuning curves in neuroscience are feature visualization in
mechanistic interpretability; lesions are ablations; connectomics is circuit
discovery; dimensionality reduction is probing's cousin; Granger causality is the
correlational ancestor of causal attribution. Each pair shares a method and a
failure mode. \emph{Hence} a result measured on one transfers as a caution to the
others: the tuning-curve trap and the probing trap are the same trap, and the danger
zone of Fig.~\ref{fig:faithfulness} is populated by both traditions at once. This
cross-disciplinary reach is the reason a benchmark built on an Atari chip speaks to
systems neuroscience, to explainable AI, to mechanistic interpretability, and to
software verification in the same breath.

\subsection*{Is a 1977 chip a fair test? The necessary-condition screen}

The obvious objection deserves a full hearing. The VCS is deterministic,
hand-coded, and small; the systems these methods are built for are learned,
distributed, stochastic, and large. Why should a score here predict anything about a
transformer? Our answer is that we do not ask it to. The benchmark is a
\emph{necessary-condition screen}, not a sufficiency predictor. A method that fails
\emph{here}---with perfect ground truth, full observability, exact interventions,
and no learning to muddy the picture---has no business being trusted on a neural
network, where every one of those advantages is gone. Passing here is necessary, not
sufficient; failing here is decisive.

\emph{In fact}, the same logic upgrades the semantic-gap claim, and this is the part
worth stating carefully. If even the \emph{strongest} causal account that a
fully observable machine permits stops short of meaning, then on a neural
network---less observable, with no clean ground truth---faithfulness is \emph{a
fortiori} not understanding. The gap we measure here is therefore a floor on the gap
there. We measure $1-S=0.56$ on an exactly recovered circuit; we should not expect
to do better where the circuit itself is uncertain.

The objection also has a constructive answer: the VCS \emph{does} exhibit the
properties that make real systems hard to interpret. Its RAM cells are aliased and
polysemantic, reused across routines, a hand-coded echo of superposition; it
computes by racing the electron beam, a distributed temporal computation; it reads a
floating bus; its registers mean different things in different contexts.
Figure~\ref{fig:representativeness} maps each VCS property to the neural-network
difficulty it does and does not capture---a property provable here is a property
suspected there. What the chip does \emph{not} capture we state without hedging:
learned features, scale, and stochasticity are absent, and the learned agent, which
folds all three back in, is Paper~5. The map is the honest accounting of where the
analogy holds and where it breaks.

\begin{figure}[t]
\centering
\includegraphics[width=\textwidth]{fig5_representativeness_map.pdf}
\caption{\textbf{The VCS$\leftrightarrow$neural-network failure-mode map --- why a
1977 chip is a fair screen.} Each row is an interpretability failure mode we
\emph{measured} on the VCS (left, blue), paired with the well-documented
neural-network counterpart it stands in for (right, vermillion). The status column
carries the necessary-condition argument: on the VCS each failure is \emph{provable}
because we hold the exact ground truth, whereas on a neural network the same failure
is only \emph{suspected}. Among the rows: vanilla saliency and integrated gradients
score $0.000$ on the discrete position regime (the shattered-gradient analogue);
Granger causality reaches faithfulness $0.136$ with a high false-edge rate
(the spurious-edge analogue); the single-unit lesion view is faithful ($0.987$) yet
\emph{misses} the wiring that connectomics recovers ($0.41$); and the sparse
autoencoder matches every named variable yet scores $F=0.041$ (the
present-not-used and superposition analogue). \emph{Note that} passing this screen
is necessary, not sufficient: a method that fails here---with perfect ground truth,
full observability, exact interventions, and no learning---has no business being
trusted where every one of those advantages is gone. What the VCS does \emph{not}
capture (learned features, scale, stochasticity) is stated in the legend and
deferred to Paper~5. Every VCS value reads from a committed record.}\label{fig:representativeness}
\end{figure}

\subsection*{What interpretability needs}

Three recommendations follow, and we state them flatly. \emph{First}, prefer causal
over correlational evidence. The danger zone is not a hypothetical: the methods the
field reaches for first carry the highest plausibility proxy and the lowest measured
faithfulness, and a map that is provably empty on the position regime looks
\emph{more} convincing than the map that is exactly right. \emph{Second}, validate
on known systems before trusting on unknown ones---sieve the methods through a
ground-truth screen, because a method that cannot pass where the answer is computable
cannot be audited where it is not. \emph{Third}, and most important, stop treating
faithfulness as understanding. The open problem for explainable AI is not a better
saliency map; it is the semantic gap. Faithfulness is the half we now know how to
measure and, on causal methods, how to win. Meaning is the half that remains.

\subsection*{Limitations}

We have measured a gap; we have not closed it, and several limits bound exactly what
we may conclude. \emph{First}, we did not run semantic recovery. Phase~E
(documentation and design recovery) is specified but not executed, so our claim is
the \emph{bounded} one---on exact ground truth, faithfulness certifies the wiring,
and the wiring alone is not the semantics---and \emph{not} the unbounded ``no method
can recover meaning,'' which would require running the recovery and is Paper~4's
charge. \emph{Second}, the sufficiency and minimality axes are causal quantities
defined against the oracle, and the T3 semantic labels are \emph{external}, imported
from AtariARI~\cite{anand2019unsupervised} and OCAtari~\cite{delfosse2023ocatari} and
only causally verified, never discovered. We say this plainly, and we turn it into
the point: even the strongest causal account a fully known machine permits stops
short of meaning, and recovering meaning from scratch is a different problem.
\emph{Third}, a single low sufficiency could be a tuning artifact. The claim
therefore does not rest on any one number; it rests on the gap appearing where
tuning cannot reach---an already-exact effect table---and on its running in both
directions on the same axis at once. \emph{Fourth}, the plausibility axis is a
documented method-tradition proxy, not a human-subjects measurement, and we report
the danger zone at that proxy's coarse resolution rather than overclaiming it.
\emph{Fifth}, our scores hold within Paper~1's validated conformance window; outside
that horizon the substrate is not certified bit-exact, and we keep every scored claim
inside it. \emph{Finally}, the subject is one old machine with no learning. That is
the necessary-condition framing again, conceded once more: a floor, not a forecast.

\subsection*{Beyond the brain, and toward the learned agent}

We close by reaching past the present result. Because the substrate is
differentiable, it tells us which \emph{gradient} explanations are trustworthy and
where they vanish---the content path carries an honest gradient, while the discrete
index does not, so a saliency map over sprite position is empty by construction, not
by accident. That is a statement about gradients as a class, and it travels to any
system where a discrete decision sits downstream of a continuous one. We believe the
broader value of a fully known, differentiable testbed is exactly this: it lets the
community calibrate its instruments before deploying them where calibration is
impossible. The hard case remains the learned agent, where there is no disassembly,
no specification, no engineer's intent to recover, and where the semantic gap we
measured here can only widen. It is foreseeable that closing that gap---turning a
faithful account of a learned system into an understanding of what it has
learned---will be the work of the next decade of interpretability. We have not solved
it. We have made it a number, and a number is where understanding begins.

% =====================================================================
% BIBTEX-NEEDED: full entries for the keys THIS section introduces that are NOT
% yet declared in references.bib or in a sibling section's BIBTEX-NEEDED block.
% (SM to merge; do NOT edit references.bib here.)
%
% Already present / declared elsewhere, reused as-is (NO entry below):
%   jonas2017could, marr1982vision, hewitt2019designing, anand2019unsupervised,
%   delfosse2023ocatari (declared in 03_methods.tex BIBTEX-NEEDED);
%   leaderboardP2 (declared in 04_results_A.tex); faithfuldemoP2 (declared in
%   07_results_compare.tex).
%
% NEW keys introduced here (verified citations; SM to add to references.bib):
%
% @inproceedings{lindner2023tracr,
%   author    = {Lindner, David and Kram{\'a}r, J{\'a}nos and Farquhar, Sebastian and
%                Rahtz, Matthew and McGrath, Thomas and Mikulik, Vladimir},
%   title     = {Tracr: Compiled Transformers as a Laboratory for Interpretability},
%   booktitle = {Advances in Neural Information Processing Systems (NeurIPS)},
%   year      = {2023},
%   note      = {arXiv:2301.05062}
% }
%
% @inproceedings{gupta2024interpbench,
%   author    = {Gupta, Rohan and Yu, Eric and Marks, Samuel and Adams, Keith and
%                Garriga-Alonso, Adri{\`a}},
%   title     = {{InterpBench}: Semi-Synthetic Transformers for Evaluating Mechanistic
%                Interpretability Techniques},
%   booktitle = {Advances in Neural Information Processing Systems (NeurIPS),
%                Datasets and Benchmarks Track},
%   year      = {2024},
%   note      = {arXiv:2407.14494}
% }
%
% @misc{yang2019bim,
%   author       = {Yang, Mengjiao and Kim, Been},
%   title        = {{BIM}: Towards Quantitative Evaluation of Interpretability Methods
%                   with Ground Truth},
%   year         = {2019},
%   howpublished = {arXiv preprint},
%   note         = {arXiv:1907.09701}
% }
%
% @article{lazebnik2002can,
%   author  = {Lazebnik, Yuri},
%   title   = {Can a Biologist Fix a Radio? --- Or, What I Learned While Studying
%              Apoptosis},
%   journal = {Cancer Cell},
%   volume  = {2},
%   number  = {3},
%   pages   = {179--182},
%   year    = {2002}
% }
%
% @misc{ieee610,
%   author       = {{IEEE}},
%   title        = {{IEEE} Standard Glossary of Software Engineering Terminology
%                   (IEEE Std 610.12-1990)},
%   year         = {1990},
%   howpublished = {IEEE, New York},
%   note         = {Defines ``software'' as computer programs, procedures, and
%                   possibly associated documentation and data.}
% }
% =====================================================================
